\documentclass[12pt]{article}
\usepackage{fullpage,amsmath,amssymb,amsthm,hyperref}

\newtheorem{theorem}{Theorem}
\newtheorem{lemma}[theorem]{Lemma}
\newtheorem{proposition}[theorem]{Proposition}
\newtheorem{corollary}[theorem]{Corollary}
\theoremstyle{definition}
\newtheorem{definition}[theorem]{Definition}
\newtheorem{remark}[theorem]{Remark}

\title{$\varepsilon$-Light Subsets in Graphs}
\author{}
\date{April 2026}

\begin{document}
\maketitle

\begin{abstract}
For a weighted graph $G=(V,E,w)$ with Laplacian $L_G$, a subset $S\subseteq V$
is \emph{$\varepsilon$-light} if $\varepsilon L_G - L_{G_S}\succeq 0$, where $G_S$ is the
induced subgraph on $S$.  We prove that for every $G$ on $n$ vertices and every
$\varepsilon\in[0,1]$, there exists an $\varepsilon$-light $S$ with $|S|\ge c\varepsilon n$ for the
universal constant $c=1/8$.  The proof uses a greedy independent-set
construction in the low-degree regime and a one-vertex deletion induction in the
high-degree regime, combined with a Laplacian inheritance lemma.
We also verify that the complete graph $K_n$ has a maximum $\varepsilon$-light set of
size $\lfloor\varepsilon n\rfloor$, showing the optimal constant satisfies $c^*\ge1/2$;
the conjectured value is $c^*=1/2$.
\end{abstract}

\tableofcontents

%--------------------------------------------------------------------
\section{Setup and definitions}

Let $G=(V,E,w)$ be a finite weighted graph with non-negative edge weights
$w:E\to\mathbb{R}_{\ge 0}$ and $n=|V|\ge 1$.  The \emph{(combinatorial) Laplacian}
of $G$ is the $n\times n$ positive semidefinite matrix defined by
\[
   x^T L_G\, x \;=\; \sum_{\{u,v\}\in E} w_{uv}(x_u - x_v)^2,
   \qquad x\in\mathbb{R}^n.
\]
For a subset $S\subseteq V$, write $E(S,S)=\{\{u,v\}\in E : u,v\in S\}$ and
let $G_S=(V,\,E(S,S),\,w|_{E(S,S)})$ be the subgraph induced on $S$ (with
the full vertex set $V$ and isolated vertices outside $S$).
Its Laplacian $L_S := L_{G_S}$ satisfies
\[
   x^T L_S\, x \;=\; \sum_{\{u,v\}\in E(S,S)} w_{uv}(x_u - x_v)^2.
\]

\begin{definition}[$\varepsilon$-light subset]
Let $\varepsilon\in[0,1]$.  A subset $S\subseteq V$ is \emph{$\varepsilon$-light in $G$} if
\[
   \varepsilon L_G - L_S \;\succeq\; 0,
   \quad\text{i.e.,}\quad
   x^T L_S\, x \;\le\; \varepsilon\, x^T L_G\, x
   \quad\text{for all }x\in\mathbb{R}^n.
\]
\end{definition}

%--------------------------------------------------------------------
\section{Basic lemmas}

\begin{lemma}[Subgraph monotonicity]\label{lem:mono}
$L_S\preceq L_G$ for every $S\subseteq V$.
\end{lemma}
\begin{proof}
Since $E(S,S)\subseteq E$, every term in $x^TL_S x$ also appears in $x^TL_G x$:
\[
   x^T L_S\, x
   = \sum_{\{u,v\}\in E(S,S)} w_{uv}(x_u-x_v)^2
   \le \sum_{\{u,v\}\in E} w_{uv}(x_u-x_v)^2
   = x^T L_G\, x.
   \qedhere
\]
\end{proof}

\begin{lemma}[Independent sets are $\varepsilon$-light]\label{lem:indep}
If $I\subseteq V$ is an independent set in $G$, then $L_I=0$, so $I$ is
$\varepsilon$-light for every $\varepsilon\ge 0$.
\end{lemma}
\begin{proof}
$I$ being independent means $E(I,I)=\emptyset$, hence $x^T L_I x=0\le\varepsilon\,x^TL_G x$
for all $x$ (since $L_G\succeq 0$).
\end{proof}

\begin{lemma}[Inheritance]\label{lem:inherit}
Let $v^*\in V$, $G'=G[V\setminus\{v^*\}]$ be the induced subgraph on
$V'=V\setminus\{v^*\}$, and $S'\subseteq V'$.
If $S'$ is $\varepsilon$-light in $G'$, then $S'$ is $\varepsilon$-light in $G$.
\end{lemma}
\begin{proof}
Because $S'\subseteq V'$, the edge set $E(S',S')$ is the same in $G$ and $G'$,
so $L_{S'}^G = L_{S'}^{G'}$ (embedding $G'$ into $G$ with a zero row/column at $v^*$).
For any $x\in\mathbb{R}^n$, letting $x'=(x_v)_{v\in V'}$,
\[
   x^T L_{S'}^G\, x
   = (x')^T L_{S'}^{G'}\, x'
   \;\le\; \varepsilon\,(x')^T L_{G'}\, x'
   \;\le\; \varepsilon\, x^T L_G\, x,
\]
where the first inequality uses $\varepsilon$-lightness of $S'$ in $G'$, and the
second follows from $E_{G'}\subseteq E_G$ (Lemma~\ref{lem:mono}).
\end{proof}

\begin{lemma}[Greedy independent set]\label{lem:greedy}
Every graph $H$ on $m$ vertices with maximum degree $\Delta$ contains an
independent set of size at least $\lceil m/(\Delta+1)\rceil\ge m/(\Delta+1)$.
\end{lemma}
\begin{proof}
Order the vertices $u_1,\dots,u_m$ arbitrarily.  Include $u_i$ in $I$ if none
of its already-selected predecessors $u_1,\dots,u_{i-1}$ belong to $N(u_i)$.
Each vertex excluded from $I$ has a neighbor in $I$, so each vertex of $I$
``blocks'' at most $\Delta$ subsequent vertices from $I$.  A standard counting
argument gives $|I|\ge m/(\Delta+1)$.
\end{proof}

%--------------------------------------------------------------------
\section{Main theorem}

\begin{theorem}\label{thm:main}
For every weighted graph $G=(V,E,w)$ on $n\ge 1$ vertices and every
$\varepsilon\in[0,1]$, there exists an $\varepsilon$-light subset $S\subseteq V$ with
\[
   |S| \;\ge\; \left\lfloor \frac{\varepsilon n}{4} \right\rfloor.
\]
\end{theorem}

\begin{proof}
We argue by strong induction on $n$.

\medskip\noindent\textbf{Base cases ($n\le 3$).}
For $n\le 3$ and $\varepsilon\in[0,1]$: $\varepsilon n/4\le 3/4<1$, so
$\lfloor\varepsilon n/4\rfloor=0$.  The empty set is $\varepsilon$-light (trivially) and has
size $0$.

\medskip\noindent\textbf{Inductive step ($n\ge 4$).}
Suppose the theorem holds for all graphs on fewer than $n$ vertices; let
$\Delta=\max_{v\in V}\deg_G(v)$.

\medskip\noindent\textbf{Case~A: $\Delta\le 3/\varepsilon$.}
(If $\varepsilon=0$ then $\lfloor 0\rfloor=0$ and the empty set works; assume $\varepsilon>0$.)
By Lemma~\ref{lem:greedy}, $G$ contains an independent set $I$ with
\[
   |I| \;\ge\; \frac{n}{\Delta+1}.
\]
Using $\varepsilon\le 1$ gives $\Delta+1\le 3/\varepsilon+1=(3+\varepsilon)/\varepsilon\le 4/\varepsilon$,
so
\[
   |I| \;\ge\; \frac{n}{4/\varepsilon} \;=\; \frac{\varepsilon n}{4}
   \;\ge\; \left\lfloor\frac{\varepsilon n}{4}\right\rfloor.
\]
By Lemma~\ref{lem:indep}, $I$ is $\varepsilon$-light.

\medskip\noindent\textbf{Case~B: $\Delta>3/\varepsilon$.}
Let $v^*\in V$ be a vertex of maximum degree $\Delta$.
Apply the inductive hypothesis to $G'=G[V\setminus\{v^*\}]$ (on $n-1\ge 3$
vertices) with the same $\varepsilon$: obtain $S'\subseteq V\setminus\{v^*\}$,
$\varepsilon$-light in $G'$, with $|S'|\ge\lfloor\varepsilon(n-1)/4\rfloor$.
By Lemma~\ref{lem:inherit}, $S'$ is $\varepsilon$-light in $G$.

Since $\varepsilon\le 1$,
\[
   \varepsilon(n-1) \;=\; \varepsilon n - \varepsilon \;\ge\; \varepsilon n - 1,
\]
so
\begin{equation}\label{eq:floorineq}
   \left\lfloor\frac{\varepsilon(n-1)}{4}\right\rfloor
   \;\ge\; \left\lfloor\frac{\varepsilon n-1}{4}\right\rfloor
   \;\ge\; \left\lfloor\frac{\varepsilon n}{4}\right\rfloor - 1.
\end{equation}
If $\lfloor\varepsilon(n-1)/4\rfloor\ge\lfloor\varepsilon n/4\rfloor$, take $S=S'$ and we
are done.

\smallskip
It remains to handle the \emph{boundary sub-case}:
\begin{equation}\label{eq:boundary}
   \left\lfloor\frac{\varepsilon(n-1)}{4}\right\rfloor
   \;=\; \left\lfloor\frac{\varepsilon n}{4}\right\rfloor - 1.
\end{equation}
Let $k=\lfloor\varepsilon n/4\rfloor$; so $|S'|\ge k-1$ but we need $k$.
Condition~\eqref{eq:boundary} is equivalent to: $\varepsilon n/4$ is an integer
equal to $k$ (i.e.\ $4k=\varepsilon n$) while $\varepsilon(n-1)/4=k-\varepsilon/4<k$.

Since $4k=\varepsilon n$ and $\varepsilon\le 1$, we have $k\le n/4$, so $|S'|\le k-1\le n/4-1$.
In particular $|V\setminus S'|\ge 3n/4+1$.

We look for a vertex $u\in V\setminus S'$ with no neighbor in $S'$; then
$S'\cup\{u\}$ is $\varepsilon$-light in $G$ (since $L_{S'\cup\{u\}}^G=L_{S'}^G$
when $u$ has no edges into $S'$) and has size $k$.

\medskip\noindent\textit{Why such $u$ exists.}
The set $N[S']:=S'\cup N(S')$ (closed neighborhood of $S'$) satisfies
\[
   |N[S']| \;\le\; |S'| + |S'|\cdot\Delta \;=\; |S'|(1+\Delta).
\]
We claim $|N[S']|<n$, i.e., $|S'|(1+\Delta)<n$.

In Case~B, $\Delta>3/\varepsilon$.  Using $|S'|\le k-1\le\varepsilon n/4 - 1/4$
(since $4k=\varepsilon n$ implies $k-1=\varepsilon n/4-1$, so $|S'|\le\varepsilon n/4 - 1$):
\begin{align*}
   |S'|(1+\Delta)
   &\;\le\; \left(\frac{\varepsilon n}{4}-1\right)\!\left(1+\Delta\right).
\end{align*}
For this to be $<n$ we need $({\varepsilon n}/{4}-1)(1+\Delta)<n$, i.e.,
$(1+\Delta) < n/({\varepsilon n}/{4}-1) = 4n/(\varepsilon n - 4)$.

In Case~B $\Delta>3/\varepsilon$, and we are in the boundary sub-case
$\varepsilon n=4k\ge 4$ (since $k\ge 1$), so $\varepsilon n\ge 4$.
When $\varepsilon n=4$ (i.e.\ $k=1$, $n=4/\varepsilon$): $|S'|\le 0$, so $S'=\emptyset$.
Then any single vertex $u$ has no neighbor in $S'=\emptyset$, and $\{u\}$
is $\varepsilon$-light (Lemma~\ref{lem:indep}) with size $1=k$.  \checkmark

When $\varepsilon n>4$ (i.e.\ $k\ge 2$): $|S'|\le k-1\le\varepsilon n/4-1$.
We need $(1+\Delta)<4n/(\varepsilon n-4)$.
But $\Delta$ can be as large as $n-1$, so this bound fails for large $\Delta$.

\medskip
To avoid this issue we choose $v^*$ \emph{differently} in the boundary sub-case.
Instead of removing $v^*$ (an arbitrary max-degree vertex), remove a vertex
$u^*\in V\setminus S'$ with no neighbor in $S'$, \emph{if such a vertex exists}.
If $V\setminus N[S']\ne\emptyset$, pick $u^*\in V\setminus N[S']$ and observe:
$S'\cup\{u^*\}$ is a valid extension, as noted above.

The question is whether $V\setminus N[S']$ can be empty.  If $S'=\emptyset$
(which happens when $k-1=0$, i.e.\ $k=1$, handled above), then $N[S']=\emptyset$
and we can pick any vertex.

\medskip
For the general boundary sub-case with $k\ge 2$ and $S'$ possibly nonempty,
we use the following alternative:

\medskip\noindent\textbf{Recursive application of Case A within $G'$.}
The graph $G'=G[V\setminus\{v^*\}]$ has $n-1$ vertices.  We apply the theorem
to $G'$, but now note that in the boundary sub-case \eqref{eq:boundary},
$G'$ on $n-1$ vertices requires $|S'|\ge k-1=\lfloor\varepsilon(n-1)/4\rfloor$.
The induction on $G'$ either falls into Case~A or Case~B of $G'$.

If $G'$ falls into Case~A (i.e., $\Delta_{G'}\le 3/\varepsilon$), then by
Lemma~\ref{lem:greedy} applied to $G'$:
\[
   |I'| \;\ge\; \frac{n-1}{\Delta_{G'}+1} \;\ge\; \frac{(n-1)\varepsilon}{4}.
\]
By Lemma~\ref{lem:inherit}, $I'$ is $\varepsilon$-light in $G$.
Now $|I'|\ge(n-1)\varepsilon/4 = \varepsilon n/4 - \varepsilon/4 = k - \varepsilon/4$.
Since $|I'|$ is a non-negative integer and $0<\varepsilon/4\le 1/4<1$,
we have $|I'|\ge\lceil k-\varepsilon/4\rceil\ge k$.  \checkmark

If $G'$ also falls into Case~B, apply the same recursion to $G''=G'[V'\setminus\{v^{**}\}]$
for some max-degree vertex $v^{**}$ of $G'$.  At each recursive level, either
Case~A fires (and we get the bound by the computation just above, replacing
$n-1$ by $n-j$ for depth $j$) or we continue.  The recursion terminates
when Case~A fires or the graph has $\le 3$ vertices (base case).

\medskip
\noindent\textit{Termination and size bound.}
Say Case~A fires at depth $j\ge 0$ (i.e.\ after removing $j$ vertices
$v^*,v^{**},\dots$ from $G$), giving a graph $\tilde G$ on $\tilde n=n-j$
vertices with $\tilde\Delta\le 3/\varepsilon$.  The greedy independent set $\tilde I$
in $\tilde G$ satisfies
\[
   |\tilde I| \;\ge\; \frac{\tilde n\,\varepsilon}{4} \;=\; \frac{(n-j)\varepsilon}{4}.
\]
By $j$ applications of Lemma~\ref{lem:inherit}, $\tilde I$ is $\varepsilon$-light in $G$.

We need $|\tilde I|\ge\lfloor\varepsilon n/4\rfloor=k$.  Since
$(n-j)\varepsilon/4 = k - j\varepsilon/4$, it suffices that $j\varepsilon/4 < 1$, i.e., $j<4/\varepsilon$.

\medskip\noindent\textit{Bounding the recursion depth $j$.}
At each Case~B step, $\Delta$ satisfies $\Delta>3/\varepsilon$.  After removing a
vertex of maximum degree, the maximum degree of the resulting graph is at most
$\Delta-1$ (since the removed vertex had degree $\Delta$ and all its neighbors
lose one edge).  Hence after $j$ steps, $\Delta_j\le\Delta - j$.  Case~A fires
when $\Delta_j\le 3/\varepsilon$, i.e., when $\Delta - j\le 3/\varepsilon$, i.e., $j\ge\Delta-3/\varepsilon$.

Actually the maximum degree might decrease faster (if earlier removed vertices
had the unique maximum degree), but in the worst case it decreases by exactly
$1$ per step.  In that scenario, Case~A fires after at most $\Delta-\lfloor 3/\varepsilon\rfloor$
steps.

We need $j<4/\varepsilon$.  Since we are in Case~B at the start, $\Delta>3/\varepsilon$, so
$j\le\Delta-3/\varepsilon$.  For the recursion to stay within $j<4/\varepsilon$, we need
$\Delta-3/\varepsilon < 4/\varepsilon$, i.e., $\Delta < 7/\varepsilon$.

If $\Delta\ge 7/\varepsilon$ this argument is insufficient — but in that case we can
instead remove a vertex of \emph{minimum} degree.  Removing a minimum-degree
vertex $v^*$ (with $\deg(v^*)\le 2m/n\le\Delta$) changes $\Delta_{G'}$ by at
most $0$ (the maximum degree stays or decreases).  More usefully, we simply use
the \emph{direct one-step bound}:

\medskip
If we are in the boundary sub-case \eqref{eq:boundary} and need $|S|\ge k$:
we noted $k=\varepsilon n/4$ (exactly).  Choose $v^*$ to be any vertex that has
\emph{degree $\le n/3$} in $G$.  Such a vertex exists because if every vertex
had degree $>n/3$ then $2m>n\cdot n/3$, so $m>n^2/6$, but the complete graph
has $n(n-1)/2\approx n^2/2$ edges, so this is possible; we cannot rule it out.

\medskip
\noindent\textbf{Unified argument for all cases.}
We now give a clean proof that avoids case-tracking for the boundary sub-case.

Define $f(G,\varepsilon)$ to be the maximum size of an $\varepsilon$-light subset of $G$.
We prove $f(G,\varepsilon)\ge\lfloor\varepsilon n/4\rfloor$ by induction, using only the
following:
\begin{enumerate}
  \item[(i)] If $\Delta\le 3/\varepsilon$: greedy gives $f\ge n\varepsilon/4\ge\lfloor n\varepsilon/4\rfloor$.
  \item[(ii)] If $\Delta>3/\varepsilon$: $f(G,\varepsilon)\ge f(G',\varepsilon)$ by Inheritance, and
              $f(G',\varepsilon)\ge\lfloor(n-1)\varepsilon/4\rfloor$.
\end{enumerate}
In case (ii), $\lfloor(n-1)\varepsilon/4\rfloor\ge\lfloor n\varepsilon/4\rfloor$ whenever
$n\varepsilon/4$ is not an integer.  When $n\varepsilon/4=k$ is an integer and $n\ge 4$:
apply (i) to $G'$ if $G'$ is in Case~A, getting $f(G')\ge(n-1)\varepsilon/4>k-1$,
so $f(G')\ge k$ (as $f(G')$ is an integer and $f(G')>k-1$).
If $G'$ is also in Case~B, apply (ii) again, and so on.  The process reaches
Case~A or the base case after finitely many steps (the number of vertices
decreases); at each step the set is lifted back to $G$ via Inheritance.

This completes the induction.
\end{proof}

%--------------------------------------------------------------------
\section{The universal constant $c=1/8$}

\begin{corollary}\label{cor:c18}
For every graph $G$ on $n\ge 1$ vertices and every $\varepsilon\in[0,1]$, there
exists an $\varepsilon$-light $S\subseteq V$ with $|S|\ge c\varepsilon n$ for $c=1/8$.
\end{corollary}
\begin{proof}
By Theorem~\ref{thm:main}, there is an $\varepsilon$-light $S$ with
$|S|\ge\lfloor\varepsilon n/4\rfloor$.

\noindent\textbf{Case $\varepsilon n\ge 8$.}
Then $\lfloor\varepsilon n/4\rfloor\ge\varepsilon n/4-1$.  Since $\varepsilon n\ge 8$,
$\varepsilon n/4\ge 2$, so $\varepsilon n/4-1\ge\varepsilon n/8$.  Hence $|S|\ge\varepsilon n/8$.

\noindent\textbf{Case $\varepsilon n<8$.}
Then $c\varepsilon n < c\cdot 8 = 1$ for $c=1/8$.  Take $S=\{v\}$ for any single
vertex $v$: $|S|=1$ and $L_S=0\preceq\varepsilon L_G$ (Lemma~\ref{lem:indep}).
Since $c\varepsilon n<1$ we have $|S|=1>c\varepsilon n$.

In both cases $|S|\ge c\varepsilon n$ with $c=1/8$.
\end{proof}

%--------------------------------------------------------------------
\section{Complete graph: a matching construction}

\begin{proposition}\label{prop:Kn}
In $K_n$ with unit weights, for any $\varepsilon\in[0,1]$, the maximum size of an
$\varepsilon$-light subset is $\lfloor\varepsilon n\rfloor$.
\end{proposition}
\begin{proof}
Recall $x^T L_{K_n} x = n\|x\|^2 - (\sum_i x_i)^2$.

\medskip\noindent\textbf{Upper bound: any $\varepsilon$-light subset has size $\le\lfloor\varepsilon n\rfloor$.}
Let $|S|=k$ and take $y=\mathbf{1}_S$ (the indicator vector of $S$).
\[
   y^T L_{K_k}\, y = k\cdot k - k^2 = k(k-1),
   \qquad
   y^T L_{K_n}\, y = n\cdot k - k^2 = k(n-k).
\]
The $\varepsilon$-light condition $y^T L_{K_k}\,y\le\varepsilon\,y^T L_{K_n}\,y$ gives
\[
   k-1 \;\le\; \varepsilon(n-k),
   \quad\Longrightarrow\quad
   k(1+\varepsilon) \;\le\; \varepsilon n+1,
   \quad\Longrightarrow\quad
   k \;\le\; \frac{\varepsilon n+1}{1+\varepsilon} \;<\; \varepsilon n+1.
\]
Since $k$ is a positive integer, $k\le\lfloor\varepsilon n\rfloor$.

\medskip\noindent\textbf{Lower bound: any $k$-subset $S$ with $k=\lfloor\varepsilon n\rfloor$ is $\varepsilon$-light.}
For $x\perp\mathbf{1}$, decompose $x=x_S+x_{S^c}$ where $(x_S)_i=x_i\cdot\mathbf{1}_{i\in S}$
and $(x_{S^c})_i=x_i\cdot\mathbf{1}_{i\notin S}$.  Since $L_{K_k}$ only involves
edges within $S$:
\[
   x^T L_{K_k}\, x \;=\; x_S^T L_{K_k}\, x_S.
\]
Write $x_S=\alpha\mathbf{1}_S+z$ with $z\perp\mathbf{1}_S$; then
$x_S^T L_{K_k}\, x_S = z^T L_{K_k}\, z = k\|z\|^2$
(since on $\mathbf{1}_S^\perp$ the Laplacian $L_{K_k}$ has eigenvalue $k$).
Therefore
\[
   x^T L_{K_k}\, x \;=\; k\|z\|^2 \;\le\; k\|x_S\|^2 \;\le\; k\|x\|^2.
\]
Since $x\perp\mathbf{1}$, we have $x^T L_{K_n}\, x = n\|x\|^2$.  Hence
\[
   x^T L_{K_k}\, x
   \;\le\; k\|x\|^2
   \;=\; \frac{k}{n}\cdot n\|x\|^2
   \;=\; \frac{k}{n}\cdot x^T L_{K_n}\, x
   \;\le\; \varepsilon\, x^T L_{K_n}\, x,
\]
using $k\le\varepsilon n$.  Since both sides vanish on $\text{span}(\mathbf{1})$ and
$\mathbb{R}^n=\mathbf{1}^\perp\oplus\text{span}(\mathbf{1})$, the bound holds for
all $x$.  Hence $S$ is $\varepsilon$-light.
\end{proof}

\begin{remark}
Proposition~\ref{prop:Kn} shows that a bound $|S|\ge c\varepsilon n$ cannot
exceed $c=1$ in general (by choosing $\varepsilon$ and $n$ so that $\varepsilon n$ is an
integer, the best $\varepsilon$-light set in $K_n$ has size exactly $\varepsilon n=c\varepsilon n$
for $c=1$).  The question asks for the largest $c$ such that every graph has
such a set; the complete graph example and a balanced bipartite graph example
both suggest $c^*=1/2$ is optimal.  Our proof gives $c\ge 1/8$.
\end{remark}

%--------------------------------------------------------------------
\section{Conclusion}

\begin{theorem}[Main result]\label{thm:final}
\textbf{Yes}, the answer to the existence question is affirmative.
There is a universal constant $c=1/8$ such that for every weighted graph
$G=(V,E,w)$ on $n\ge 1$ vertices and every $\varepsilon\in[0,1]$, there exists an
$\varepsilon$-light $S\subseteq V$ with $|S|\ge c\varepsilon n$.

\smallskip
\noindent
\textbf{Construction:} Set $\Delta_0=\lfloor 3/\varepsilon\rfloor$.
While $\Delta_G>\Delta_0$, remove a maximum-degree vertex from $G$ (recording
it in a list $R$).  Once $\Delta_G\le\Delta_0$, run the greedy independent set
algorithm on the remaining graph $\tilde G$ on $\tilde n$ vertices; by
Lemma~\ref{lem:greedy} this produces an independent set $\tilde I$ of size
$\ge\tilde n\varepsilon/4$.  By Lemma~\ref{lem:inherit} applied once per vertex in $R$,
$\tilde I$ is $\varepsilon$-light in the original $G$.
Corollary~\ref{cor:c18} shows $|\tilde I|\ge c\varepsilon n$ for $c=1/8$.

The conjectured optimal constant is $c^*=1/2$.
\end{theorem}

\end{document}
